\chapter*{ТЕОРЕТИЧЕСКАЯ ЧАСТЬ}
\addcontentsline{toc}{chapter}{ТЕОРЕТИЧЕСКАЯ ЧАСТЬ}

\section*{Шифр Виженера}

\textbf{Шифр Виженера} --- полиалфавитный шифр подстановки, использующий ключевое слово для циклического сдвига символов. Каждый символ ключа задаёт сдвиг в соответствующей позиции текста.

\subsection*{Принцип работы}

\begin{enumerate}
    \item Ключ повторяется циклически над всем текстом.
    \item Для каждой буквы текста вычисляется сдвиг, соответствующий букве ключа.
    \item Буква заменяется на букву алфавита, сдвинутую на это количество позиций.
\end{enumerate}

\textbf{Формула шифрования:}
$$C_i = (P_i + K_{i \bmod m}) \bmod n$$

где $P_i$ --- позиция символа открытого текста, $K_j$ --- позиция символа ключа, $m$ --- длина ключа, $n$ --- размер алфавита.

\textbf{Формула расшифрования:}
$$P_i = (C_i - K_{i \bmod m} + n) \bmod n$$

\section*{Метод Казиски}

\textbf{Метод Казиски} --- метод криптоанализа, основанный на поиске повторяющихся n-грамм в шифротексте. Расстояния между их появлениями позволяют предположить длину ключа через вычисление НОД найденных расстояний.

\subsection*{Алгоритм}

\begin{enumerate}
    \item Найти все повторяющиеся последовательности из 3+ символов в шифротексте.
    \item Вычислить расстояния между позициями каждого повторения.
    \item Найти НОД всех расстояний --- это вероятная длина ключа или её делитель.
\end{enumerate}

\subsection*{Ограничения метода}

Метод не всегда даёт точный результат, потому что:
\begin{itemize}
    \item повторения могут происходить случайно;
    \item текст может быть коротким;
    \item длинный ключ даёт мало повторов;
    \item некоторые языковые структуры порождают естественные повторения.
\end{itemize}

\section*{Индекс совпадений}

\textbf{Индекс совпадений (IC)} --- статистическая мера, показывающая вероятность того, что две случайно выбранные буквы текста окажутся одинаковыми:
$$IC = \frac{\sum_{i=0}^{n-1} f_i(f_i - 1)}{N(N-1)}$$

где $f_i$ --- частота появления $i$-й буквы, $N$ --- длина текста.

Для русского языка $IC \approx 0.0553$, для английского $IC \approx 0.0667$. Для случайного текста $IC \approx 1/n$.

При правильной длине ключа подпоследовательности имеют IC близкий к естественному языку.

\section*{Метод Кирхгофа (частотный анализ)}

При известной длине ключа $L$ шифротекст можно разбить на $L$ групп, каждая из которых зашифрована шифром Цезаря. Для каждой группы проводится частотный анализ с использованием известных статистик языка.

\subsection*{Алгоритм взлома}

\begin{enumerate}
    \item Разбить шифротекст на $L$ подпоследовательностей: первая содержит символы с позициями 0, L, 2L..., вторая --- 1, L+1, 2L+1 и т.д.
    \item Для каждой подпоследовательности подсчитать частоты букв.
    \item Для каждого возможного сдвига (0..n-1) сравнить частотный профиль с эталонным методом $\chi^2$.
    \item Выбрать сдвиг с минимальным значением $\chi^2$.
    \item Восстановить букву ключа по найденному сдвигу.
\end{enumerate}

\subsection*{Критерий хи-квадрат}

$$\chi^2 = \sum_{i=0}^{n-1} \frac{(O_i - E_i)^2}{E_i}$$

где $O_i$ --- наблюдаемая частота, $E_i$ --- ожидаемая частота для языка.

Чем меньше значение $\chi^2$, тем лучше совпадение распределений.

\section*{Эталонные частоты русского алфавита}

\begin{center}
\begin{tabular}{|c|c||c|c||c|c|}
\hline
Буква & Частота & Буква & Частота & Буква & Частота \\
\hline
О & 0.1097 & Е & 0.0845 & А & 0.0801 \\
И & 0.0735 & Н & 0.0670 & Т & 0.0626 \\
С & 0.0547 & Р & 0.0473 & В & 0.0454 \\
Л & 0.0440 & К & 0.0349 & М & 0.0321 \\
Д & 0.0298 & П & 0.0281 & У & 0.0262 \\
\hline
\end{tabular}
\end{center}
